\subsubsection{结果分析}\label{sec:q3-4}

\textbf{（1）实验设置。}策略的效果在本地复原的模拟器上检验。这个模拟器按题目附件 2 的接口说明和演练程序的生成规则复原：源在半径 $1770$ m 的圆内按面积均匀分布，数量在 $10$ 到 $16$ 之间等概率，频道从 $20$ 个中不放回抽取，接收半径在 $1000$ 到 $1500$ m 之间均匀抽取；测向误差是在 $150$ m 网格上平滑变化的空间误差场，先量化到 $0.01^\circ$ 再限制在 $\pm1^\circ$ 内，同一位置、同一频道的误差固定，重复检测不能消除；计时按题目规则，每段移动的时间取整到微秒。它不是官方程序本身，没有和官方程序做过逐输入输出的对照，因此这里的数字不能当作正式测试成绩，只用来在相同场景上比较不同策略。

比较五种策略，都由同一套几何模块搭成，区别只在规则：
\begin{itemize}
\item \textbf{本文方案}：即算法\ref{alg:q3-search}，含提前光学清除；
\item \textbf{无提前光学}：去掉（5）的规则，区域必须缩到 $19.5$ m 以内才清除，其余相同；
\item \textbf{原点先扫}：在无提前光学的基础上，开局先在原点把 $20$ 个频道扫一遍；
\item \textbf{联合规划}：在无提前光学的基础上，排路线时把已知源的位置也当作可能的扫描点，若能替代某个站就剪掉那个站；
\item \textbf{早期版本}：候选探测点只在圆心两侧横向取 $6$ 个、路线只用最近邻加 2-opt 的更早一版。
\end{itemize}
场景共 $2000$ 个，每个场景五种策略从同一初始状态出发，源位置、接收半径和误差场完全相同，共 $10000$ 次运行；场景在策略冻结之后才生成，没有用来调参数。评价指标是每个场景的总耗时除以该场景的源数，再对场景等权平均；总耗时包括清除最后一个源之后为确认没有遗漏而做的扫描。策略之间的差用同场景配对的 $10000$ 次自助法给出 $95\%$ 区间。另用 $50$ 个场景、每种策略各跑两遍、单进程顺序执行，测每局的计算时间。这五种策略是从更多候选中按此前 $184$ 局官方演练的成绩挑出来的，但那些演练没有让各策略跑相同的场景，不能用来排名，所以本文的比较以下面 $2000$ 个同场景的结果为准。

\textbf{（2）主要结果。}$10000$ 次运行全部完整清除，共 $26006$ 个源实例，没有失败或超时的场景。各策略的平均定位清除时间见表\ref{tab:q3-main}，相对无提前光学方案的同场景差值见表\ref{tab:q3-paired}。

\Needspace{46mm}
\begin{center}
\begin{minipage}{\linewidth}
\centering\small
\captionof{table}{五种策略在 $2000$ 个相同场景上的平均定位清除时间（单位：s/源）}
\label{tab:q3-main}
\vspace{4pt}
\begin{tabular}{@{}lrcrrrr@{}}
\toprule
策略 & 均值 & 均值 95\% 区间 & 中位数 & P95 & 全清场景 & 计算时间 / s/局 \\
\midrule
本文方案 & $239.26$ & {[}237.85, 240.68{]} & $236.60$ & $296.29$ & $2000/2000$ & $0.057$ \\
无提前光学 & $240.89$ & {[}239.45, 242.36{]} & $238.07$ & $297.68$ & $2000/2000$ & $0.049$ \\
联合规划 & $240.92$ & {[}239.46, 242.38{]} & $238.17$ & $297.76$ & $2000/2000$ & $0.066$ \\
原点先扫 & $241.38$ & {[}239.87, 242.87{]} & $238.81$ & $300.39$ & $2000/2000$ & $0.056$ \\
早期版本 & $243.85$ & {[}242.34, 245.36{]} & $241.62$ & $303.70$ & $2000/2000$ & $0.125$ \\
\bottomrule
\end{tabular}
\end{minipage}
\end{center}

\Needspace{40mm}
\begin{center}
\begin{minipage}{\linewidth}
\centering\small
\captionof{table}{相对无提前光学方案的同场景差值（负值表示更快，单位：s/源）}
\label{tab:q3-paired}
\vspace{4pt}
\begin{tabular}{@{}lrcr@{}}
\toprule
策略 & 平均差值 & 差值 95\% 区间 & 更快 / 相同 / 更慢的场景数 \\
\midrule
本文方案 & $-1.631$ & {[}$-1.715$, $-1.546${]} & $1736\ /\ 0\ /\ 264$ \\
联合规划 & $+0.032$ & {[}$-0.019$, $+0.084${]} & $86\ /\ 1821\ /\ 93$ \\
原点先扫 & $+0.493$ & {[}$+0.034$, $+0.946${]} & $955\ /\ 0\ /\ 1045$ \\
早期版本 & $+2.958$ & {[}$+2.477$, $+3.433${]} & $760\ /\ 0\ /\ 1240$ \\
\bottomrule
\end{tabular}
\end{minipage}
\end{center}

本文方案的均值最低，比无提前光学方案每源平均省 $1.63$ s，区间不含零；$2000$ 个场景中 $1736$ 个更快、$264$ 个更慢，没有一个相同，说明提前光学几乎在每一局都改变了动作序列。另外三种对照各自说明一件事：联合规划和无提前光学在 $1821$ 个场景上动作完全相同，其余场景互有快慢，平均差值的区间跨零，也就是把已知源位置当扫描点来剪站，在这个问题上基本没有机会用上；原点先扫平均慢 $0.49$ s/源，区间刚好不含零，胜负场景接近一半对一半，说明开局那 $119$ s 的检测费用大体上被早发现源省下的路程抵消，但抵消得不够；早期版本平均慢 $2.96$ s/源，是候选探测点和路线搜索两处改进的合计收益。

\textbf{（3）时间省在哪里。}把本文方案和无提前光学方案每局的动作分开看：平均检测次数由 $123.8$ 次降到 $118.0$ 次，少了 $5.7$ 次，约 $29$ s；平均移动距离由 $11358$ m 增到 $11410$ m，多走 $52$ m，约 $10$ s；提前尝试在 $2000$ 局里共失败 $1616$ 次，平均每局 $0.8$ 次，约 $2.4$ s。三项合计与每局约 $21$ s 的实际节省相符（每局平均 $13.0$ 个源）。也就是说，提前光学省下的是本来要在近距离上再做的一两次测向；代价是失败的尝试和去重心多走的路，两者都比省下的检测便宜。这 $1616$ 次失败全部按 $3$ s 计入总耗时，没有从结果里去掉。

\textbf{（4）按源数分组。}表\ref{tab:q3-count} 和图\ref{fig:q3-strata} 按场景的真实源数（只在事后分组时使用）列出结果。源越多，每源平均时间越短，因为扫站的固定开销被更多的源分摊；源数为 $16$ 的场景还能在发现第 $16$ 个源后停止扫站。本文方案在七个分组里都比无提前光学快，差值稳定在 $1.5$ 到 $1.8$ s/源之间，与源数无关；原点先扫在源少的分组里更慢、源多的分组里略快，因为源多时开局多扫一遍更容易早发现几个源。

\Needspace{44mm}
\begin{center}
\begin{minipage}{\linewidth}
\centering\small
\captionof{table}{按源数分组的平均定位清除时间（单位：s/源）}
\label{tab:q3-count}
\vspace{4pt}
\begin{tabular}{@{}lrrrrrrr@{}}
\toprule
源数 & 10 & 11 & 12 & 13 & 14 & 15 & 16 \\
\midrule
场景数 & 295 & 268 & 299 & 291 & 259 & 297 & 291 \\
本文方案 & $286.95$ & $267.45$ & $250.87$ & $235.47$ & $223.15$ & $213.51$ & $197.43$ \\
无提前光学 & $288.48$ & $269.11$ & $252.52$ & $237.03$ & $224.78$ & $215.14$ & $199.20$ \\
联合规划 & $288.52$ & $269.28$ & $252.56$ & $237.07$ & $224.70$ & $215.19$ & $199.17$ \\
原点先扫 & $290.46$ & $270.70$ & $252.73$ & $238.64$ & $224.48$ & $214.17$ & $198.54$ \\
早期版本 & $293.16$ & $273.16$ & $255.03$ & $241.29$ & $226.90$ & $216.67$ & $200.77$ \\
\bottomrule
\end{tabular}
\end{minipage}
\end{center}

\begin{center}
\begin{minipage}{\linewidth}
\centering
\QThreeStrataFigure
\captionof{figure}{按源数分组的结果。(a) 四种策略的平均定位清除时间都随源数增加而下降。(b) 相对无提前光学方案的同场景差值，本文方案在每个分组都在零线以下，联合规划贴着零线。}
\label{fig:q3-strata}
\end{minipage}
\end{center}

\textbf{（5）个例。}平均改善不等于每个场景都改善。以逐例记录中的几个场景为例：编号 f392164b 的 $11$ 源场景，本文方案 $239.83$ s/源，无提前光学 $249.36$ s/源，省了 $9.5$ s/源，是提前尝试连续命中的情形；编号 2baf85fc 的 $11$ 源场景，本文方案 $274.28$ s/源，比无提前光学的 $274.10$ 慢 $0.18$ s/源，是一两次尝试落空、又多走了一段路的情形。慢的 $264$ 个场景里，差值多数在 $1$ s/源以内，这与每源最多两次、每次 $3$ s 的费用上限一致；去重心测向引起的路线变化才是更慢场景的主要来源，这一部分没有上限，但从分布看很少超过几秒。

\textbf{（6）计算开销。}表\ref{tab:q3-main} 最后一列是单进程顺序复跑 $50$ 个场景得到的每局计算时间，本文方案 $0.057$ s，比无提前光学多 $0.008$ s，多出的部分是面积样本点的生成和占比计算；早期版本最慢，因为它的路线搜索没有编译。计算时间都远小于机器狗一局几千秒的动作时间。$10000$ 次运行的每一个动作都按记录重放核对过，反馈和逐段计时一致；每个场景都检查了退出条件，源数少于 $16$ 时所有未发现频道在每个已扫描站都无信号且未去的站为空，源数为 $16$ 时有 $16$ 次真实的清除成功。

\textbf{（7）适用条件。}以上结果只对本文冻结的参数和本地模拟器的场景生成方式成立。提前光学的门槛 $0.4$、$80$ m 和 $2$ 次是开发阶段定下的，没有按这 $2000$ 个场景重新调过；如果官方演练的误差特性与此差别较大，提前尝试的命中率可能变化，但它的费用上限不变，几何部分的保证也不受影响：无论尝试成功与否，站扫完就不会漏源，判定能清除时一定能清除。官方演练和正式测试的清除比例与平均定位清除时间，要以实际日志为准。
